\section{Uniform fillings with finitely many simplex types}
\label{sec:fillings}

We now prove the chain-level part of Theorem~\ref{thm:torus}.  Its
purpose is to obtain one iterate of $f$ and one finite set of allowable
simplex types that fill every cycle of a specified dimension and input
bound.  The companion~\cite[Lemma~2.5 and Proposition~2.6]{OpenAIOvergroups}
uses the same homological
device.  It requires neither uniformly short fillings
nor an algorithm for finding them.

\subsection{The homogeneous complex and its bounds}

Let $\mathcal E(U)$ be the simplicial set whose $q$-simplices are all
tuples $(u_0,\ldots,u_q)$ of elements of $U$.  Faces delete entries, and
degeneracies repeat entries.  Write $C_*(U)$ for its normalized integral
chain complex.  Thus its basis in degree $q$ consists of tuples with
no two consecutive entries equal; a tuple with consecutive equal
entries represents zero.  The differential is zero in degree zero,
and, for $q\ge1$, it is
\[
 \partial(u_0,\ldots,u_q)
   =\sum_{i=0}^q(-1)^i(u_0,\ldots,\widehat{u_i},\ldots,u_q).
\]
The augmentation $\epsilon:C_0(U)\to\Z$ sends every vertex to one.
Prepending the identity vertex contracts the augmented complex as a
complex of abelian groups.  Simultaneous left translation gives a free
$U$-action on the normalized basis, so the augmented complex is also
a free resolution of the trivial left $\Z U$-module.  The contraction
need not be $U$-equivariant.

A \emph{bound} is a finite symmetric subset $A\subseteq U$ containing
$1$.  Let $\mathcal E_A(U)$ be the simplicial subset of tuples with
\[
 u_i^{-1}u_j\in A\qquad\text{for all }i,j,
\]
and denote its normalized chains by $C_*^A(U)$.  The condition is
preserved by deleting or repeating vertices, so this is a chain
subcomplex.  A reduced zero-cycle means an element of the kernel of
$\epsilon$; in positive degree, ``reduced cycle'' means an ordinary
cycle.

In each fixed degree, $C_*^A(U)$ has only finitely many basis elements
up to left translation.  Indeed, translating the first vertex to $1$
puts every remaining vertex in $A$.  These orbits are the
\emph{simplex types} of bound $A$.  Notice the distinction: a bound
restricts each simplex's relative vertices, not the locations of its
vertices in $U$, the number of simplices in a chain, or its integer
coefficients.

A homomorphism of groups sends a finite bound to a finite bound.
We say that a chain operator on $C_*^A(U)$ is \emph{controlled} through
specified degrees if there is a finite bound containing every simplex
in all of its output chains in those degrees.  For finitely many
operators and degrees, separate output bounds may be replaced by
their union, with inverses and $1$ adjoined.  For chains in
$\mathcal E(U\times U)$, we will also use separate finite bounds in
the two coordinates.

\subsection{A controlled additivity calculation}

Let $c:C_*(U)\to C_*(U)$ denote the chain map of the constant vertex
map $u\mapsto1$.  Thus $c_0$ sends every vertex to $(1)$, and $c_q=0$
for $q>0$.  The next lemma is an ordinary chain calculation: its
homotopies are not required to commute with the $U$-action.

\Needspace{20\baselineskip}
\begin{lemma}\label{lem:controlled-additivity}
Let $n\ge1$, let $A$ be a bound, and let $f:U\to U$ be an
endomorphism.  Suppose that, for some $k\ge0$, the induced chain map
$F=(f^k)_*:C_*^A(U)\to C_*(U)$ has homomorphisms
\[
 H_q:C_q^A(U)\longrightarrow C_{q+1}(U)
       \qquad(0\le q<n)
\]
whose output chains have a common finite bound and which satisfy
\[
 \partial H_q+H_{q-1}\partial=F_q-c_q
       \qquad(0\le q<n),
\]
where $H_{-1}=0$ and the differential on $C_0$ is zero in this
unaugmented equation.  Let $h_1,h_2:U\to U$ be homomorphisms with
elementwise commuting images, and put $h(u)=h_1(u)h_2(u)$.
Then there is a finite bound $B$ such that every $n$-cycle
$z\in C_n^A(U)$ satisfies
\begin{equation}\label{eq:controlled-additivity}
 [(hf^k)_*z]=[(h_1f^k)_*z]+[(h_2f^k)_*z]
       \quad\text{in }H_n(C_*^B(U)).
\end{equation}
One bound $B$ may be chosen for any fixed finite collection of pairs
$(h_1,h_2)$ satisfying these hypotheses.
\end{lemma}

\begin{proof}
Extend $H$ by zero in all degrees $q\ge n$, and define the chain map
\[
 F'=F-\partial H-H\partial.
\]
Then $F'_0=c_0$ and $F'_q=0$ for $1\le q<n$.  The truncation $H_n=0$
is important: if $z$ is an $n$-cycle, it gives
\begin{equation}\label{eq:truncated-homotopy}
 F'z=Fz.
\end{equation}
Although the two maps need not agree on arbitrary $n$-chains, this
equality is exactly what will be needed.  A common finite output
bound controls $F,F'$ and $H$ in every degree used below; taking
boundaries does not enlarge a bound.

We use the classical Alexander--Whitney and shuffle comparison in the
setting of Eilenberg--Zilber~\cite{EilenbergZilber1953}.  We spell out
the chain maps and their support control: the later product-module
argument requires one finite set of simplex types across all
coordinates.  The Alexander--Whitney diagonal is
\begin{equation}\label{eq:aw-diagonal}
 D(u_0,\ldots,u_q)
   =\sum_{i=0}^q
      (u_0,\ldots,u_i)\otimes(u_i,\ldots,u_q).
\end{equation}
We use the tensor differential
\[
 \partial(a\otimes b)=\partial a\otimes b
                   +(-1)^{\deg a}a\otimes\partial b.
\]
In the boundary of \eqref{eq:aw-diagonal}, terms on the two sides of
an interior cut cancel, and the remaining terms give $D\partial$.
Thus $D$ is a chain map.  It is well defined on normalized chains:
if the input has consecutive equal vertices, every summand has a
degenerate factor.  Both tensor factors retain the input bound.

The shuffle map
\[
 S:C_*(U)\otimes C_*(U)\longrightarrow C_*(U\times U)
\]
sends $(u_0,\ldots,u_p)\otimes(v_0,\ldots,v_q)$ to the signed sum
over monotone lattice paths from $(0,0)$ to $(p,q)$.  A path contributes
the tuple obtained by replacing each of its vertices $(i,j)$ by
$(u_i,v_j)$.  Its sign is $(-1)^N$, where $N$ counts pairs of steps
in which a vertical step precedes a horizontal step.  In its
boundary, deleting a corner pairs two paths that traverse the corner
in opposite orders; these contributions cancel.  The remaining
faces are precisely the shuffle of the tensor differential, so $S$
is a chain map.  If either input is degenerate, every path has a
step with identical endpoint pairs and hence gives a degenerate
output.  The map therefore descends to normalized chains and
preserves the bound in each coordinate.

Write $\Delta$ for the chain map of the diagonal homomorphism
$U\to U\times U$.  There is a chain homotopy $R$ satisfying
\begin{equation}\label{eq:diagonal-homotopy}
 \partial R+R\partial=\Delta-SD
\end{equation}
and preserving the bound in both coordinates.  Here is a direct
construction that also verifies this control.  On vertices take
$R_0=0$, since $\Delta$ and $SD$ agree there.  Suppose $R$ has been
defined in lower degrees, and let $a=(u_0,\ldots,u_q)$ be a
normalized basis simplex of positive degree, with vertex set $V$.
The chain
\[
 \Delta a-SDa-R\partial a
\]
is a cycle: apply the previously established homotopy equation to
$\partial a$.  It lies in the full tuple complex on $V\times V$.
That complex is contracted by prepending any fixed vertex, so
prepending $(u_0,u_0)$ fills the displayed cycle.  Choose this filling
as $R(a)$ and extend linearly on the normalized basis.  Every
coordinate vertex in this construction belongs to $V$, which proves
the claimed bound.  This construction takes place on the normalized
basis; degenerate faces have already been set to zero.

Apply \eqref{eq:diagonal-homotopy} after $F$.  Naturality of $D$
for the vertex map $f^k$ gives a controlled chain homotopy from
$\Delta F$ to $S(F\otimes F)D$.  To replace $F$ by $F'$, use the
degree-one tensor operator $K$ defined on homogeneous tensors by
\begin{equation}\label{eq:tensor-homotopy-signs}
 K(a\otimes b)
     =H(a)\otimes F(b)+(-1)^{\deg a}F'(a)\otimes H(b).
\end{equation}
The tensor differential and the chain-map identities for $F,F'$
give
\begin{align*}
 \partial K+K\partial
   &=(F-F')\otimes F+F'\otimes(F-F')\\
   &=F\otimes F-F'\otimes F'.
\end{align*}
In particular the sign in the second term of
\eqref{eq:tensor-homotopy-signs} cannot be dropped.  The operator
$SKD$ is controlled separately in both coordinates: it uses only
the already controlled maps $F,F',H$ and the bound-preserving maps
$D,S$.  No equivariance of $H$ enters this calculation.
Explicitly, the composite homotopy $L=RF+SKD$ satisfies
\[
 \partial L+L\partial
       =\Delta F-S(F'\otimes F')D.
\]

Let $z$ be an $n$-cycle.  Since $F'$ vanishes in degrees
$1,\ldots,n-1$, only the terms of bidegrees $(0,n)$ and $(n,0)$
survive in $(F'\otimes F')Dz$.  The degree-zero factor always maps
to $(1)$.  Summing the endpoint terms over the whole cycle first, and
then using \eqref{eq:truncated-homotopy}, gives
\[
\begin{aligned}
 (F'\otimes F')Dz
   &=(1)\otimes F'z+F'z\otimes(1)\\
   &=(1)\otimes Fz+Fz\otimes(1).
\end{aligned}
\]
This also covers $n=1$, when there are no intermediate bidegrees.
Shuffling a vertex against a chain simply holds that coordinate
fixed.

Finally, commuting images make
\[
 m:U\times U\longrightarrow U,
       \qquad m(u,v)=h_1(u)h_2(v)
\]
a homomorphism.  If two pairs have coordinate differences in
$B_1$ and $B_2$, their image difference belongs to
$h_1(B_1)h_2(B_2)$.  Indeed, the factors from the two image subgroups
commute, allowing the two differences to be grouped.  This finite
set is symmetric and contains $1$.  Thus applying $m_*$ to all the
preceding homotopies retains finite control and proves
\eqref{eq:controlled-additivity}: the finite chain $m_*Lz$ fills
the difference of its left- and right-hand cycle representatives.
For a finite collection of pairs,
take the union of the finitely many resulting output bounds.
\end{proof}

\subsection{One iterate fills every cycle of a fixed bound}

We now combine the additivity calculation with the two finite
idempotent systems.  The exponent and the output bound in the next
statement are chosen before the input cycle, which is the uniformity
needed when passing later to products of coefficient modules.

\begin{proposition}\label{prop:uniform-fillings}
Assume the hypotheses of Theorem~\ref{thm:torus}.  For every finite
bound $A\subseteq U$ and every integer $n\ge0$, there are an integer
$k\ge0$ and a finite bound $B\subseteq U$ such that, for every
reduced integral $n$-cycle $z\in C_n^A(U)$, there is a finite chain
$y\in C_{n+1}^B(U)$ satisfying
\[
 \partial y=(f^k)_*z.
\]
The chain $y$ may depend on $z$, with no uniform restriction on its
length or coefficients.
\end{proposition}

\begin{proof}
We argue by induction on $n$.  By the finite-presentation
factorization $f=ba$ in Theorem~\ref{thm:torus}, the subgroup $b(P)$
is finitely generated and contains $f(U)$.  Choose a finite symmetric
generating set $S$ for $b(P)$, with $1$ adjoined.  For every $u\in U$,
choose a finite edge path from $1$ to $f(u)$ whose consecutive
differences belong to $S$.  Its chain $H_0(u)$ satisfies
\[
 \partial H_0(u)=(f(u))-(1).
\]
Extend $H_0$ linearly.  For a reduced zero-cycle $z$, the constant
vertex terms cancel, so $\partial H_0(z)=f_*z$.  All these paths use
the same bound $S$, even though their lengths can be arbitrarily
large.  This proves the assertion for $n=0$.

Suppose $n\ge1$ and the assertion holds in all degrees less than
$n$, for every finite input bound.  We first construct the partial
controlled homotopy in Lemma~\ref{lem:controlled-additivity} on
$C_*^A(U)$.  Start with $F=f_*$ and the vertex paths just chosen,
which give its homotopy to the constant map in degree zero.

Suppose, for some $1\le q<n$, that $F=(f^k)_*$ and controlled
operators $H_0,\ldots,H_{q-1}$ have already been constructed with
\[
 \partial H_i+H_{i-1}\partial=F_i-c_i
       \qquad(0\le i<q).
\]
For each normalized basis $q$-simplex $a\in C_q^A(U)$, the next
homotopy equation asks us to fill
\begin{equation}\label{eq:partial-homotopy-obstruction}
 F(a)-H_{q-1}(\partial a).
\end{equation}
This is a $q$-cycle.  Its boundary is zero by the preceding homotopy
equation applied to $\partial a$; the remaining constant-map term
$c_{q-1}(\partial a)$ is zero, also when $q=1$ because the
augmentation of a boundary is zero.  Moreover, all these cycles
have one common finite bound, independent of the basis simplex
$a$: the current $F$ sends $A$ to a finite bound and $H_{q-1}$ is
controlled.

Apply the induction assertion in degree $q$ to this common bound.
It supplies an integer $r\ge0$ and another common finite bound such
that $(f^r)_*$ sends every cycle
\eqref{eq:partial-homotopy-obstruction} to a boundary there.
Postcompose $F$ and the already chosen $H_i$ by $(f^r)_*$.
The existing homotopy equations remain valid, since $f^r(1)=1$.
For each basis simplex choose a finite filling of its newly mapped
cycle and use that filling as $H_q(a)$.  Linear extension defines a
controlled homomorphism $H_q$, and the next homotopy equation now
holds.

There are only finitely many degrees $q<n$, so this process ends
with one power $f^k$ and a controlled homotopy to the constant map
through degrees less than $n$.  For $n=1$ the initial vertex
homotopy already suffices.  None of these choices is required to be
equivariant or computable.  In particular, we do not choose one
filling on each orbit and then translate it.  Uniformity instead
comes from the induction assertion applying at each step to
\emph{every} cycle in a single bounded complex.

Apply Lemma~\ref{lem:controlled-additivity} to every orthogonal
relation in both $\mathcal P_2$ and $\mathcal P_3$.  There are only
finitely many relations, so there is one finite output bound $B$
for all of them; enlarge it if necessary to contain the images of
the input bound under every $h_{\ell,E}f^k$.  For a fixed $n$-cycle $z$,
write
\[
 z_{\ell,E}=[(h_{\ell,E}f^k)_*z]\in H_n(C_*^B(U)).
\]
For each $\ell\in\{2,3\}$, the commuting-image and orthogonal-sum identities,
together with Lemma~\ref{lem:controlled-additivity}, give
\[
 z_{\ell,E+P'}=z_{\ell,E}+z_{\ell,P'}
\]
for every orthogonal triple in that finite system.  Here the three maps
in each relation are evaluated on the common input $f^k(u)$.  The two
full-identity classes agree because both systems have $h_{\ell,I}=f$:
\[
 z_I:=z_{2,I}=z_{3,I}=[(f^{k+1})_*z].
\]
By Lemma~\ref{lem:idempotents}, this single class is annihilated by
both $2$ and $3$.  Hence it is zero, since $z_I=3z_I-2z_I$.
It therefore has a finite filling in $C_*^B(U)$.  The exponent
$k+1$ and the bound $B$ were chosen independently of $z$, proving
the required assertion in degree $n$ and completing the induction.
\end{proof}

Proposition~\ref{prop:uniform-fillings} is a statement about ordinary
integral chains with uniform finite simplex types.  Its homotopies
need not be equivariant, and it asserts no decision procedure for
the word problem of $U$.  The next step uses precisely this form of
control to assemble coordinatewise fillings with product
coefficients.
